%% Section 4: Theoretical Framework

\section{Theoretical Framework}
\label{sec:theory}

This section develops a formal model of enterprise behaviour under a cash
supply shock to derive testable predictions. The framework embeds heterogeneous
payment technologies into a cash-in-advance (\textsc{cia}) model of \citet{lucas1987money}.

\subsection{Environment}

\paragraph{Enterprises.}
A continuum of enterprises is indexed by $i \in [0,1]$, each located in
state $s \in \{1,\ldots,S\}$. Enterprise $i$ produces a single good with the
production function
\begin{equation}
  y_i = A_i \,\ell_i^{\alpha}, \qquad \alpha \in (0,1),
  \label{eq:prod}
\end{equation}
where $A_i > 0$ is enterprise-specific productivity and $\ell_i \geq 0$ is
labour input hired at competitive wage $w > 0$. The enterprise sells output
at price $p$ (normalised to one) and maximises profit net of a fixed
operating cost $\kappa_i > 0$:
\begin{equation}
  \pi_i \;=\; A_i \ell_i^{\alpha} - w\ell_i - \kappa_i.
  \label{eq:profit}
\end{equation}

\paragraph{Payment technology.}
Wages must be paid before output is sold, a timing friction generating a
transactions demand for a means of payment. Let $\phi_{is} \in [0,1)$ denote
the share of the wage bill that enterprise $i$ can settle through digital
channels (mobile money, USSD transfers, fintech applications). The remainder
$(1-\phi_{is})$ must be settled in physical cash. Parameterise
$\phi_{is} = \phi(F_s, \nu_i)$, where $F_s \geq 0$ is the state-level
fintech penetration index and $\nu_i$ is enterprise-specific idiosyncratic
digital access, with $\partial\phi/\partial F_s > 0$.

\paragraph{Cash-in-advance constraint.}
The enterprise must hold sufficient physical currency before hiring:
\begin{equation}
  m_{is} \;\geq\; (1-\phi_{is})\, w\,\ell_{is},
  \label{eq:cia}
\end{equation}
where $m_{is} \leq \bar{m}_{is}$ is the stock of physical cash available
to the enterprise, bounded above by local cash supply $\bar{m}_{is}$ that is
proportional to the state aggregate $M_{st}$.

\subsection{Unconstrained and Constrained Optima}

When the \textsc{cia} constraint~\eqref{eq:cia} does not bind, the enterprise
chooses the first-best labour input:
\begin{equation}
  \ell_i^{\textsc{fb}} \;=\;
    \left(\frac{\alpha A_i}{w}\right)^{1/(1-\alpha)}.
  \label{eq:lfb}
\end{equation}
The first-best profit is
$\pi_i^{\textsc{fb}} = (1-\alpha)\,A_i\,(\ell_i^{\textsc{fb}})^{\alpha} - \kappa_i$.

When the \textsc{cia} constraint binds---because local cash supply
$\bar{m}_{is}$ falls short of $(1-\phi_{is})\,w\,\ell_i^{\textsc{fb}}$---the
enterprise is rationed to:
\begin{equation}
  \ell_i^{*} \;=\;
    \frac{\bar{m}_{is}}{(1-\phi_{is})\,w},
  \label{eq:lconstrained}
\end{equation}
and constrained profit is
\begin{equation}
  \pi_i^{*} \;=\;
    A_i\!\left(\frac{\bar{m}_{is}}{(1-\phi_{is})\,w}\right)^{\!\alpha}
    - \bar{m}_{is}/(1-\phi_{is}) - \kappa_i.
  \label{eq:profit_constrained}
\end{equation}

The \textsc{cia} constraint binds if and only if
\begin{equation}
  \bar{m}_{is} \;<\; (1-\phi_{is})\,w\,\ell_i^{\textsc{fb}}.
  \label{eq:binding_condition}
\end{equation}

\subsection{The Cash Supply Shock}

Prior to the shock (October 2022), state cash supply is $M_{s,\text{pre}} = \bar{M}$
for all $s$, with $\bar{M}$ high enough that the \textsc{cia} constraint does not
bind for any enterprise. The naira crunch reduces available cash to
$M_{s,\text{crunch}} = \bar{M} - \Delta_s$, where $\Delta_s > 0$ for all $s$.
For sufficiently large $\Delta_s$, the constraint~\eqref{eq:cia} binds.

\paragraph{Exit decision.}
Enterprise $i$ exits if constrained profit falls below zero:
\begin{equation}
  \pi_i^{*} < 0
  \quad\Longleftrightarrow\quad
  \bar{m}_{is} < \underline{m}_i(\phi_{is}),
  \label{eq:exit}
\end{equation}
where $\underline{m}_i(\phi_{is})$ is an exit-threshold cash level that is
\emph{decreasing} in $\phi_{is}$, as shown in Proposition~\ref{prop:insulation}
below.

\subsection{Proposition: Digital Payment Insulation}

\begin{proposition}[Digital Payment Insulation]
\label{prop:insulation}
Suppose the \textsc{cia} constraint~\eqref{eq:cia} binds at the crunch. Then:
\begin{enumerate}[label=(\roman*)]
  \item \textbf{Employment.} Constrained employment $\ell_i^{*}$ is strictly
    increasing in the digital payment share $\phi_{is}$:
    \[
      \frac{\partial \ell_i^{*}}{\partial \phi_{is}}
        = \frac{\bar{m}_{is}}{(1-\phi_{is})^2\,w} > 0.
    \]
  \item \textbf{Survival.} The exit threshold $\underline{m}_i(\phi_{is})$
    is strictly decreasing in $\phi_{is}$: a higher digital payment share
    expands the set of cash-supply realisations under which the enterprise
    survives.
  \item \textbf{Heterogeneity.} The employment gain from a marginal increase in
    $\phi_{is}$ is strictly increasing in the cash supply contraction
    $\Delta_s$:
    \[
      \frac{\partial^2 \ell_i^{*}}{\partial \phi_{is}\,\partial \Delta_s} > 0.
    \]
\end{enumerate}
\end{proposition}

\begin{proof}
Part~(i) follows directly from differentiating~\eqref{eq:lconstrained} with
respect to $\phi_{is}$.

For part~(ii), differentiate the exit condition in~\eqref{eq:exit} implicitly.
Constrained profit $\pi_i^{*}(\bar{m}_{is}, \phi_{is})$ satisfies
$\partial\pi_i^{*}/\partial\bar{m}_{is} > 0$ (more cash raises output and
profit) and $\partial\pi_i^{*}/\partial\phi_{is} > 0$ (higher digital share
raises constrained employment and profit). Setting $\pi_i^{*} = 0$ and
applying the implicit function theorem gives
$\mathrm{d}\underline{m}_i/\mathrm{d}\phi_{is}
= -(\partial\pi_i^{*}/\partial\phi_{is})/(\partial\pi_i^{*}/\partial\bar{m}_{is}) < 0$.

Part~(iii) follows from the cross-derivative of~\eqref{eq:lconstrained}:
since $\bar{m}_{is} = \bar{M} - \Delta_s$ decreases in $\Delta_s$,
$\partial\ell_i^{*}/\partial\phi_{is} = \bar{m}_{is}/[(1-\phi_{is})^2 w]$
falls as $\Delta_s$ rises, but the \emph{level} effect of $\phi_{is}$ on
employment is larger in absolute terms when $\Delta_s$ is large, because the
binding constraint is tighter.
\end{proof}

\subsection{Testable Predictions}

The three parts of Proposition~\ref{prop:insulation} map directly to the
empirical specifications.

\paragraph{Prediction 1.} Part~(i) predicts a positive coefficient on
$\textit{Fintech}_s \times \mathbf{1}[\text{peak crunch}]$ in the enterprise
employment regression (Section~\ref{sub:nlps_results}).

\paragraph{Prediction 2.} Part~(ii) predicts a positive coefficient in the
enterprise survival linear probability model; and that the survival advantage
is larger when cash supply falls further (tested across shock phases in
Section~\ref{sub:nlps_results}).

\paragraph{Prediction 3.} Part~(iii) predicts that the insulation effect is
larger for enterprises with higher baseline cash dependence---specifically
small and rural enterprises for whom the \textsc{cia} constraint is most
likely to bind. Tested in Section~\ref{sub:heterogeneity}.

\paragraph{Prediction 4.} Aggregating across enterprises within a state, a
higher average $\phi_s = \mathbb{E}_i[\phi_{is}|s]$ reduces state-level
output loss. This is validated using the nighttime-light pre-trend test in
Section~\ref{sub:pretrendvalid}.

\subsection{Theory Discussion}

The insulation channel operates through \emph{payment medium substitution},
not credit access. Enterprise $i$ is not assumed to borrow more or have
better collateral in high-fintech states; it simply settles a larger fraction
of its wage bill without physical cash, relaxing the binding
constraint~\eqref{eq:cia}. This distinguishes the mechanism from the
financial accelerator \citep{repec:eee:macchp:1-21}, where the amplification
channel runs through collateral values and external finance premiums.

The model also predicts that insulation effects are proportional to the
\emph{magnitude} of the cash supply contraction, not its level. This motivates
the event-study design in Section~\ref{sec:design}, which traces dynamics
across the announcement, peak, recovery, and post-shock phases rather than
collapsing to a single-period DiD.

\paragraph{Partial equilibrium and the fixed-wage assumption.}
The model takes the competitive wage $w$ as exogenous and common across
states. This is appropriate if labour markets are integrated nationally
rather than at the state level, or if wages are predetermined at the start
of the period. In general equilibrium, a binding \textsc{cia} constraint
in low-fintech states could reduce local labour demand and compress wages,
partially attenuating the employment decline. The employment results in
Section~\ref{sec:results_micro} are therefore conservative estimates of
the labour-market cost of the cash supply shock in low-fintech states:
the observed employment gap understates the full wedge when falling wages
are accounted for.
